\chapter{Discussion}

\section{Performance}
The dense sentence model improves over both the fair BM25 sentence baseline and the review-level
BM25 baseline. Its MAP is close to the RIRD fine-tuned reference despite using an off-the-shelf
encoder. The result is not a universal leaderboard claim because the systems, training data and
evaluation protocols differ. The meaningful conclusion is that a simple multilingual dense model
is competitive while remaining reproducible and usable for Thai transfer.

The comparison also explains why the proposed model is not simply ``the largest model''. BM25
is competitive when reviews contain query vocabulary, but dense retrieval recovers paraphrases
and related concepts. Sentence units provide better evidence for the top-$k$ configuration, but
lexical sentence retrieval loses useful context. Top-10 mean is the practical compromise
supported by the experiments.

\section{Interpretability}
The system exposes original review sentences that contribute to retrieval. This is more
actionable than a rating score and allows users to detect disagreement, outdated information or
negative evidence. However, the deletion experiment is inconclusive and a high-scoring sentence
can contradict the query. Human support annotation is therefore essential. The system does not
generate synthetic explanations, avoiding unsupported claims that were not present in the data.

\section{Complexity and deployment}
BM25 is inexpensive and interpretable at token level. Dense retrieval requires precomputed
embeddings and model memory but supports semantic and multilingual matches. A cross-encoder is
substantially slower, so the live demo restricts it to the top 20 restaurants and 15 sentences
per restaurant. Query rewriting adds an API dependency and cost, so it is cached and used only
on the live path. Precomputed benchmark queries return immediately; arbitrary English queries
require embedding and optional Groq/CE inference.

\begin{table}[H]
\centering\small
\caption{Practical comparison of system components.}
\label{tab:complexity}
\begin{tabularx}{\textwidth}{@{}lLLL@{}}
\toprule
Component & Strength & Cost & User-facing role\\
\midrule
BM25 & Exact lexical matches; deterministic & Low CPU cost & Transparent baseline\\
Bi-encoder & Semantic and multilingual matching & Embedding storage and model memory & Main ranking model\\
LLM rewrite & Expands vague language and constraints & API latency and possible drift & Query understanding\\
Cross-encoder & Joint query--sentence matching & Slow inference & Selective negation reranking\\
\bottomrule
\end{tabularx}
\end{table}

\section{Comparison with Kaggle and Hugging Face solutions}
Many available food datasets and example systems focus on star-rating prediction, sentiment
classification or nearest-neighbour search. Our system offers a different user benefit:
query-specific restaurant ranking with inspectable review evidence. Its unique features are fair
BM25/dense comparison, top-$k$ aggregation analysis, review-volume control, sentence-level
explanations, selective negation handling and a Thai taste case study. These features make it
preferable when trust and preference matching matter more than predicting a single global rating.

\section{What improved the final system}
The final system is the result of several evidence-driven changes:
\begin{enumerate}
\item \textbf{Fair unit control:} BM25 and dense retrieval use identical review and sentence
units with identical aggregation functions.
\item \textbf{Sentence-level evidence:} dense sentence retrieval with top-10 mean makes every
recommendation traceable to concrete text.
\item \textbf{Volume control:} the 400-review experiment measures and reduces unequal review
opportunities.
\item \textbf{LLM query rewriting:} Groq rewrites vague phrases such as ``cheat meal'' into
concrete reviewer language and converts ``no waiting line'' into searchable attributes.
\item \textbf{Selective cross-encoder reranking:} B3 showed that reranking every query hurts
average performance, so the live system activates it only for explicit negation and exclusion.
\item \textbf{Evidence filtering and caching:} short generic sentences are excluded from live
explanations, embeddings are precomputed, and repeated queries are cached.
\end{enumerate}
This sequence is important: the LLM and cross-encoder improve query handling and the demo
experience, but they do not replace the measured dense baseline. The final architecture keeps
the strongest validated ranking model as its default and adds complexity only where error
analysis indicates a need.

\section{Limitations and threats to validity}
RIRD contains only 50 restaurants and English queries. Wongnai lacks restaurant IDs, so Thai
results are sentence-level only. External model and Groq behaviour can change, and rewriting
may introduce unsupported attributes. Review text may be biased or stale. The cross-encoder was
not trained on this domain, and the deletion test is limited by the pooled evidence design.
Future work should use larger multilingual reviewed-item benchmarks, domain-adapted rerankers and
human-centred explanation studies.
